
% Ejercicio 3
\ejercicio{
    Se considera la aplicación proyectiva 
    \[
    g : \mathbb{P}^2_{\mathbb{C}} \longrightarrow \mathbb{P}^2_{\mathbb{C}}
    \]
    que tiene por matriz, respecto del sistema de referencia proyectivo canónico,
    \[
    \begin{pmatrix}
    0 & 1 & 2 \\
    0 & 1 & 0 \\
    0 & 0 & 2
    \end{pmatrix}.
    \]
    
    Halla centro e imagen de $g$.

}{

    La aplicación lineal asociada es
    \[
    \hat{g} : \mathbb{C}^3 \longrightarrow \mathbb{C}^3, \quad M_{B_c,B_c}(\hat{g}) = \begin{pmatrix}
    0 & 1 & 2 \\
    0 & 1 & 0 \\
    0 & 0 & 2
    \end{pmatrix}.
    \]
    El centro de $g$ viene dado por
    \[
    Z(g) = \mathbb{P}(\ker(\hat{g})) 
    \]
    \[
        \ker(\hat{g}) = L[(1,0,0)] \qquad Im(\hat{g}) = L[(1,1,0),(2,0,2)]
    \]
    \[
        Z(g) = \mathbb{P}(\ker(\hat{g})) = \{(1:0:0)\}
    \]
    \[
        Im(g) = \mathbb{P}(Im(\hat{g})) = V((1:1:0),(2:0:2)) : x_0 - x_1 + x_2 = 0
    \]

}

% Ejercicio 6
\ejercicio{

Prueba que la homografía $h$ de la recta $x_0 - x_1 + x_2 = 0$ sobre la recta 
$x_1 - x_2 = 0$ (rectas de $\mathbb{P}^2_{\mathbb{R}}$) que transforma los puntos 
\[
A = (1 : 1 : 0), \quad 
B = (1 : 2 : 1), \quad 
C = (1 : 0 : -1)
\]
en 
\[
A' = (1 : 0 : 0), \quad 
B' = (1 : 1 : 1), \quad 
C' = (1 : -1 : -1),
\]
respectivamente, es una perspectividad. 

\medskip

Halla la matriz y escribe las ecuaciones de $h$ respecto de las referencias 
\[
\mathcal{R} = \{A, B; C\} \quad \text{y} \quad 
\mathcal{R}' = \{A', B'; C'\}.
\]

}{

    $h$ perspectividad $\iff$ el punto de intersección de las rectas es un punto fijo.

    \[
    \mathfrak{R} = \{(1:1:0),(1:2:1);(1:0:-1)\}
    \]
    Veamos que es una referencia
    \[
        \exists \lambda, \lambda_0, \lambda_1 \in  \mathbb{R} : \lambda(1,0,-1) = \lambda_0(1,1,0) + \lambda_1(1,2,1)
    \]
    \[
        \implies \begin{cases}
            \lambda_0+\lambda_1 = \lambda \\
            \lambda_0 + 2\lambda_1 = 0 \\
            \lambda_1 = -\lambda
        \end{cases} \implies \begin{cases}
            \lambda_0 = 2\lambda \\
            \lambda_1 = -\lambda
        \end{cases} (\lambda = 1)
    \]
    \[
        \mathcal{B} = \{(2,2,0),(-1,-2,-1)\} \text{ base asociada a } \mathfrak{R}
    \]
    \[
        \mathfrak{R}' = \{(1:0:0),(1:1:1);(1:-1:-1)\}
    \]
    Veamos que es una referencia
    \[
        \exists \mu, \mu_0, \mu_1 \in  \mathbb{R} : \mu(1,-1,-1) = \mu_0(1,0,0) + \mu_1(1,1,1)
    \]
    \[
        \implies \begin{cases}
            \mu_0+\mu_1 = \mu \\
            \mu_1 = -\mu
        \end{cases} \implies \begin{cases}
            \mu_0 = 2\mu \\
            \mu_1 = -\mu
        \end{cases} (\mu = 1)
    \]
    \[
        \mathcal{B}' = \{(2,0,0),(-1,-1,-1)\} \text{ base asociada a } \mathfrak{R}'
    \]
    \[
        M_{\mathfrak{R},\mathfrak{R}'}(h) = \begin{pmatrix}
            1 & 0 \\
            0 & 1
        \end{pmatrix}
    \]

    \[
        r \cap r' : \begin{cases}
            x_0 = 0 \\
            x_1 - x_2 = 0
        \end{cases} = \{(0:1:1)\}
    \]
    \[
        (0,1,1) = \lambda(2,2,0) + \beta(-1,-2,-1)
    \]
    \[
        (0,1,1) = \lambda'(2,0,0) + \beta'(-1,-1,-1)
    \]
    como en ambas referencias tiene coordenadas $(1:2)$ y el cambio de la referencia $\mathfrak{R}$ a $\mathfrak{R}'$ es la identidad, el punto es fijo y por tanto $h$ es una perspectividad.

}